\begin{frame}{Dataset Instrumental: Coordenadas y Amplitudes}
    \begin{block}{Tabla 1: Datos de entrada para la Función de Error}
        \centering
        \scriptsize
        \begin{tabular}{@{}lcccc@{}}
            \toprule
            \textbf{Sensor} & \textbf{$x_i$} & \textbf{$y_i$} & \textbf{$z_i$} & \textbf{$Az_i$ observado} \\ \midrule
            S1  & -4.0 & -3.0 &  0.2 & 0.005219 \\
            S2  & -2.5 &  2.8 & -0.1 & 0.009811 \\
            S3  &  0.0 &  4.0 &  0.1 & 0.012014 \\
            S4  &  3.5 &  3.0 & -0.2 & 0.024570 \\
            S5  &  4.5 & -1.5 &  0.0 & 0.080580 \\
            S6  &  2.0 & -4.0 &  0.3 & 0.071733 \\
            S7  & -1.0 & -4.5 & -0.1 & 0.027704 \\
            S8  & -4.5 &  1.0 &  0.2 & 0.003556 \\
            \textbf{\textcolor{blue}{S9}}  & \textbf{\textcolor{blue}{0.5}} & \textbf{\textcolor{blue}{0.5}} & \textbf{\textcolor{blue}{0.0}} & \textbf{\textcolor{blue}{0.852685}} \\
            S10 &  2.8 &  0.8 & -0.3 & 0.380182 \\
            S11 & -1.8 & -1.2 &  0.1 & 0.119785 \\
            S12 &  1.0 &  3.2 &  0.2 & 0.037284 \\ \bottomrule
        \end{tabular}
    \end{block}
    \vspace{0.2cm}
    \begin{center}
        \tiny \textcolor{gray}{*Nota: Los datos de S9 resaltados son los utilizados en el ejemplo numérico de la siguiente diapositiva.}
    \end{center}
\end{frame}